%  Manuscript for Future Internet (MDPI). Build: sh tools/manuscript/build_all.sh PAPER_DIR (or latexmk -pdf main.tex after the supplement).
%  Numeric results come from generated/*.tex: macros_R.tex (analysis/analyze.py, \R...), macros_V.tex (analysis/analyze_v3.py,
%  \V...), macros_W.tex (analysis/analyze_v4.py, \W...), macros_X.tex (tools/v7/analysis_v7.py, \X...), macros_Y.tex
%  (tools/v9/analysis_v9.py, \Y...), imported under these names by tools/manuscript/import_generated.py; do not type numbers by hand.
%  Supplementary tables, figures and sections are referenced as \Sref{label}; generated/supp_xref.tex is written from
%  supplementary.aux by tools/manuscript/xref_maps.py.
\documentclass[futureinternet,article,submit,moreauthors]{Definitions/mdpi}
\usepackage[ruled,linesnumbered,vlined]{algorithm2e}
\usepackage{tabularx,booktabs,multirow,makecell}
\input{generated/macros_R.tex}
\input{generated/macros_V.tex}
\input{generated/macros_W.tex}
\input{generated/macros_X.tex}
\input{generated/macros_Y.tex}
\input{generated/supp_xref.tex}
\makeatletter
\newcommand{\Sref}[1]{\@ifundefined{suppref@#1}{\textbf{S??}\PackageWarning{suppxref}{Undefined supplementary label #1}}{\@nameuse{suppref@#1}}}
\makeatother
\newcommand{\sys}{PBTI}
\newcommand{\Sobs}{\textsf{obs}}
\newcommand{\Srel}{\textsf{rel}}
\newcommand{\pubsub}{publish/\allowbreak{}subscribe}
\newcommand{\Bhist}{B5\textsubscript{60/10}}
\newcommand{\Bev}{B5\textsubscript{ev}}
\newcommand{\Rhist}{PBTI-R\textsubscript{60/10}}
\newcommand{\Rev}{PBTI-R\textsubscript{ev}}

%=================================================================
\firstpage{1}
\makeatletter
\setcounter{page}{\@firstpage}
\makeatother
\pubvolume{1}
\issuenum{1}
\articlenumber{0}
\pubyear{2026}
\copyrightyear{2026}
\datereceived{ }
\daterevised{ }
\dateaccepted{ }
\datepublished{ }

%=================================================================
\Title{Provenance-Bound Trust-Aware Ingestion for Industrial IoT Edge-to-Cloud Data Pipelines: Per-Identity Conflict Containment and Its Security--Availability Trade-offs}
\newcommand{\orcidauthorA}{0009-0008-3618-192X} % Rajneesh Shukla
\newcommand{\orcidauthorB}{0009-0005-8312-7844} % Punit Gupta
\Author{Rajneesh Shukla $^{1,}$*\orcidA{} and Punit Gupta $^{2}$\orcidB{}}
\AuthorNames{Rajneesh Shukla and Punit Gupta}
\address{%
$^{1}$ \quad Independent Researcher, Houston, TX 77001, USA; rajneeshshukla@ieee.org\\
$^{2}$ \quad Independent Researcher, Mountain View, CA 94041, USA; punitgupta@ieee.org}
\corres{Correspondence: rajneeshshukla@ieee.org}

%=================================================================
\abstract{Authenticated industrial telemetry can contain conflicting versions of the same observation. We investigate how per-identity versus feed-level conflict containment affects correctness, availability and forged-data exposure in brokered edge-to-cloud pipelines. Provenance-Bound Trust-Aware Ingestion (PBTI) integrates signed origin and relay evidence, policy-controlled admission, per-identity conflict handling and auditable recovery; a controlled feed-level comparator shares the other mechanisms. In single-host simulations, the comparator matched PBTI's final forged admissions in all \VPOneIdentical\ main scenarios and final losses in \VOTwoBFiveIdentical; the remaining difference reflected reviewer-enabled recovery. With one-shot 60~s review, feed-level containment delayed genuine data by an additional \WAoneEfiveBFivehistvsPBTIWdshort~s per observation under stolen keys while reducing forged-data exposure in a mixed workload; with matched event-driven reviewers, the primary outcomes were identical or equivalent within predefined margins in four workloads. Under equivocation with suppressed messages, feed-level holds retained fewer forgeries; selective holds removed them in amendment-informed cells at a cost in genuine data but failed their registered grid criteria. No evaluated rule contained a consistent impersonator holding the relevant keys. PBTI achieved \RTpRatioBTwo\ times the in-process throughput and \RWireRatioBTwo\ times the message size of a signature-free Zero-Trust baseline. The findings characterize conditional security--availability trade-offs, not general superiority or deployment validation.}

\keyword{Industrial Internet of Things; Zero Trust; data provenance; trustworthy data ingestion; edge-to-cloud; conflict containment; equivocation; replay protection; data lineage; MQTT}

%=================================================================
\begin{document}

\section{Introduction}\label{sec:intro}

Industrial organizations move process data from operational technology (OT) through edge gateways and message brokers into enterprise data platforms~\cite{iic_iira_v110}. Zero Trust removes implicit trust based on network location~\cite{nist_sp800_207}, and industrial Zero-Trust work concentrates on remote access, segmentation, lightweight authentication and broker-side access control~\cite{federici2023zt,zanasi2024flexzt,zhang2026lztaka,james2025mqttzt}: these mechanisms decide \emph{who may connect and what they may access}.

A data platform must also decide \emph{which data object becomes authoritative, why, and whether that decision can be re-checked later}. In brokered \pubsub\ deployments session identity ends at the broker: MQTT delivery guarantees hold per hop~\cite{oasis_mqtt5}, so a subscriber learns neither which device produced a message nor which edge workload relayed it unless that evidence travels with the message, and OPC UA PubSub's replay protection depends on optional configuration~\cite{opcua_part2}. Retransmission after reconnects, late backlogs and clock faults complicate the decision: a legitimate retransmission and a malicious replay are byte-identical, and a firmware fault and an attacker holding a stolen key can both produce two authenticated observations with the same identity.

Once two authenticated versions of one event exist, the pipeline must contain the conflict: it can withhold only the affected event identities and keep the rest of the device's feed available, or withhold the feed from the fork point until a reviewer resolves it, as signed-log systems do when they ignore or expose a forked feed~\cite{tarr2019ssb,haeberlen2007peerreview}. The first choice limits availability loss but leaves a forged version authoritative whenever no contradiction arrives; the second withholds forged and genuine data alike for as long as review takes. We found no comparison of the two granularities for IIoT telemetry under matched admission, credentials and review (Section~\ref{sec:related}); this describes our search, not a priority claim.

This study makes that comparison with \emph{Provenance-Bound Trust-Aware Ingestion} (\sys), an ingestion architecture that carries origin, relay, authorization and history evidence with each observation across a broker, and a feed-level comparator, B5, that shares every \sys\ mechanism except the conflict response. We ask:
\begin{itemize}
\item[] \textbf{RQ1} When cryptographically authenticated IIoT telemetry contains conflicting event versions, what are the correctness, availability, contamination-duration and recovery consequences of per-identity versus feed-level conflict containment under matched admission and review assumptions?
\item[] \textbf{RQ2} Under which attack and delivery-fault conditions is either granularity, or a selective-hold combination of the two, preferable, and under which conditions do all of them fail?
\item[] \textbf{RQ3} What processing, message-size, memory and storage costs does object-level provenance binding add on the evaluated host, and which deployment properties remain unmeasured?
\end{itemize}

The contributions differ in kind: (C1, \emph{architectural integration and implementation}) an ingestion architecture for observations crossing an untrusted broker that combines an origin statement, a relay countersignature naming an authorized edge workload, versioned path-bound admission, credential lifecycle with retroactive taint, signed review and Merkle-checkpointed evidence, with a reference implementation (Section~\ref{sec:arch}); (C2, \emph{elementary formal refinement}) Rule L, a per-identity conflict rule with a locality lemma that restates the scope of its update rule: without intervening lifecycle transitions, $k$ admitted statements change the class and authority of at most $2k$ event identities (Proposition~\ref{prop:locality}); (C3, \emph{experimental characterization}) a comparison with a feed-level comparator, registered in advance in the repository (self-recorded), that receives the same keys, policy, evidence and, in matched arms, the same reviewer and schedule, so that it can only withhold more and expose less and the experiments measure by how much and when (Sections~\ref{sec:methods} and~\ref{sec:results}); (C4, \emph{empirical finding}) a time-dependent characterization of the trade-off (exposure, withheld time, deadline availability, recovery); and (C5, \emph{empirical finding}) the failure conditions: full impersonation with both keys, suppressed contradictions, the costs of selective holds and reviewed release under compromised credentials. Every building block has a published analogue (Table~\ref{tab:related}); we claim no new cryptographic primitive, no superiority of either strategy, no deployment validation and no evidence that an authenticated measurement is physically correct.

\section{Related Work and Research Gap}\label{sec:related}

\textbf{Access-level Zero Trust.} NIST SP 800-207 places per-request access decisions at decision and enforcement points~\cite{nist_sp800_207}. Industrial work applies it to remote maintenance~\cite{federici2023zt}, decentralized overlays with certificate-based identity and revocation~\cite{zanasi2024flexzt}, lightweight key agreement~\cite{zhang2026lztaka}, broker-side controls~\cite{james2025mqttzt} and device credentials split into technical and operational credentials~\cite{lazzaro2026vciot}; all decide at the point of access, and data-pillar guidance labels and tags data objects without addressing their provenance~\cite{nsa2024datapillar}.

\textbf{End-to-end protection through brokers.} ENTRUST (per-topic keys and per-publisher counters that detect duplicated or dropped messages)~\cite{dahlmanns2021entrust}, MQTT-I (sequence numbers and signed, chained Merkle roots)~\cite{buccafurri2024mqtti}, SMQTT (key-server-managed authenticity and replay protection)~\cite{yousaf2026smqtt} and OSCORE-IoT~\cite{dighriri2026oscoreiot} protect \pubsub\ end to end, AUPS is an authenticated \pubsub\ system for the IoT~\cite{rizzardi2016aups}, and blockchain anchoring adds accountability~\cite{lazzaro2024mqtt}. In the IETF ACE \pubsub\ profile for CoAP, an Internet-Draft, publishers countersign their group-encrypted objects and subscribers keep a per-publisher replay window that discards a second object with an already seen sequence number; the draft describes no reporting of such conflicts~\cite{ietf_ace_coap_pubsub}. The relay statement of \sys\ binds the material of an RFC~9338 countersignature~\cite{rfc9338} and adds the workload identity and path authorization that RFC~9338 leaves to the application. Sparkplug numbers messages modulo 256 and covers security only in non-normative guidance~\cite{sparkplug3}.

\textbf{Signed logs, transparency and provenance.} Signed provenance chains~\cite{hasan2009tos}, tamper-evident logs~\cite{crosby2009tamper} and Certificate Transparency~\cite{rfc9162} provide append-only, tamper-evident evidence; the IETF SCITT architecture registers signed statements under a registration policy, with receipts and registration checks that auditors can reproduce~\cite{rfc9943}, but defines no per-event delivery-history semantics and leaves revocation strategies for compromised keys out of scope. IoT provenance schemes use physical unclonable functions or wireless fingerprints~\cite{aman2017provenance,aman2019provenance}, blockchains~\cite{javaid2018blockpro,sigwart2019bcprov}, device-signed hash-linked readouts~\cite{saito2025authenticity} or signed records with a key-compromise history~\cite{tsai2026actuation}; Sigwart et al. leave the correctness of records at entry, the admission question addressed here, to future work~\cite{sigwart2019bcprov}. C2PA judges signer revocation against a trusted time-stamp~\cite{c2pa25}; our compromise time corresponds to the X.509 invalidity date~\cite{rfc5280}. Lineage is traced in stream processing~\cite{palyvos2019genealog}, for machine-learning artefacts~\cite{sukumaran2026provenance} and in W3C PROV~\cite{w3c_prov_dm}, and intrusion-recovery systems repair dependent state from externally identified attacks or retroactive patches~\cite{goel2005taser,chandra2011warp}.

\textbf{Responses to forks.} SUNDR's fork consistency makes a server's concealment of updates detectable once clients compare views~\cite{li2004sundr}; Secure Scuttlebutt, whose feeds share the structure of our observation statement (author key, sequence, previous hash, Ed25519 over canonical JSON~\cite{ssb_protocol_guide}), ignores a feed from the fork onward~\cite{tarr2019ssb}; PeerReview turns equivocation into verifiable evidence that auditors replay~\cite{haeberlen2007peerreview}; the Global Data Plane names holes and branches in signed edge-data streams and accepts them as a partial order~\cite{mor2019gdp}; Byzantine-fault-tolerant CRDTs accept equivocating updates as concurrent operations and still converge~\cite{kleppmann2022bftcrdt}; and the 2P-BFT-Log preprint lets correct replicas agree on the earliest fork and a proof of it, after which the forked log cannot be extended~\cite{lavoie2023twopbftlog}. They address detection and convergence, not the availability of individual events of a feed.

\textbf{Edge-to-cloud data engineering.} WedgeChain lets a cloud node certify digests of client-signed entries committed at an untrusted edge~\cite{nawab2021wedgechain}. Kafka's idempotent producer de-duplicates retries by producer identifier, epoch and sequence~\cite{kip98}, and CloudEvents consumers may treat events with the same source and identifier as duplicates~\cite{cloudevents102}; both make a conflicting second version indistinguishable from a retry. IIoT pipeline studies evaluate store-and-forward recovery without security~\cite{matias2026resilient}, a counter reset after a reboot enables replay in LoRaWAN~\cite{yang2018lorawan} (hence the boot epoch), and trust scores rate data~\cite{lim2010trust} or data sources~\cite{gigli2025zonia} rather than bind admission to identity.

\textbf{Research gap.} Table~\ref{tab:related} separates what prior work provides from what this study integrates, changes and evaluates. SCITT already provides policy-gated registration with reproducible checks, so \sys\ is a telemetry-specific integration, not a new transparency primitive. What we did not find in the works examined (Table~\Sref{tab:S-coverage} and those above) is an evaluation of the \emph{conflict-response granularity} itself for IIoT telemetry, with a feed-level comparator that shares all other mechanisms, matched review schedules and time-dependent outcomes; that evaluation is the main addition.

\begin{table}[!htbp]
\caption{Elements of \sys\ with the kind of contribution (in italics), what the closest prior work already provides, what \sys\ integrates or changes, and what this study evaluates.\label{tab:related}}
\begin{adjustwidth}{-\extralength}{0cm}
\footnotesize
\renewcommand\tabularxcolumn[1]{p{#1}}% top-aligned like the first column
\begin{tabularx}{\fulllength}{>{\raggedright\arraybackslash}p{3.0cm}>{\raggedright\arraybackslash}X>{\raggedright\arraybackslash}X>{\raggedright\arraybackslash}X}
\toprule
\textbf{Element} (\emph{type}) & \textbf{Closest prior work: what it provides} & \textbf{What \sys\ integrates or changes} & \textbf{Evaluated here}\\
\midrule
Origin statement with (device, boot epoch, sequence) and predecessor hash (\emph{application}) & signed feeds with sequence and previous hash~\cite{tarr2019ssb,ssb_protocol_guide}; producer epochs and sequences~\cite{kip98,sparkplug3}; hash-linked readouts~\cite{saito2025authenticity} & epoch namespace per attesting gateway for legacy devices & replay, reboot and sequence-fault scenarios (E6--E12)\\
Relay statement binding a workload identity to one signed observation (\emph{integration}) & COSE countersignatures~\cite{rfc9338}; publisher-signed objects through an untrusted broker~\cite{ietf_ace_coap_pubsub}; workload identity~\cite{spiffe_spec} & countersignature by the authorized relaying workload, checked against path policy & stolen device keys injected at the broker (E5), ablation\\
Versioned path-bound admission and credential lifecycle (\emph{integration}) & SCITT registration policy~\cite{rfc9943}; X.509 invalidity date~\cite{rfc5280}; revocation versus time-stamp~\cite{c2pa25} & IIoT policy schema (asset, schema, relay scope, origin type); taint by compromise time & unauthorized producers (E4), revocation (E14), Rego parity\\
Per-identity conflict rule, Rule L (\emph{formal refinement; experimental characterization}) & fork consistency~\cite{li2004sundr}, detection~\cite{tarr2019ssb} and evidence~\cite{haeberlen2007peerreview}; branches~\cite{mor2019gdp}; fork convergence~\cite{lavoie2023twopbftlog,kleppmann2022bftcrdt}; first-writer-wins and replay windows~\cite{kip98,cloudevents102,ietf_ace_coap_pubsub} & the same fork evidence (conflicting versions, contradicting \textsf{prev} links) contests only the affected identities, not the feed; locality bound $2k$ (Proposition~\ref{prop:locality}) & per-identity vs feed-level containment under matched review; bounded checks\\
Evidence, decision replay, lineage, taint and review (\emph{integration}) & transparency logs and receipts~\cite{rfc9162,rfc9943}; log replay by auditors~\cite{haeberlen2007peerreview}; intrusion recovery~\cite{goel2005taser,chandra2011warp}; lineage~\cite{palyvos2019genealog,w3c_prov_dm} & revocation-triggered taint scoped to the origin key and selected by commit time (it over-approximates); signed, replayable release decisions & insider tampering (E16), taint (E14), reviewed release under revoked keys\\
Time-dependent evaluation of containment (\emph{evaluation methodology}) & final-state outcome counts in prior pipeline evaluations & exposure, withheld time, deadline availability, censoring, alerts per (identity, content) pair & all confirmatory comparisons\\
\bottomrule
\end{tabularx}
\end{adjustwidth}
\end{table}

\section{System and Threat Model}\label{sec:model}

Figure~\ref{fig:arch} shows the data path. Signing-capable devices sign each observation; for legacy devices the polling edge workload is the origin signer (gateway attestation). Each edge workload has a short-lived credential and a durable store-and-forward buffer, and messages cross an MQTT broker that keeps mutual TLS~\cite{rfc9846} and access-control lists as defence in depth but is not trusted for data-object integrity. The cloud-side verifier is a policy enforcement point for data objects: it consults an identity authority (certificate validation, rotation, revocation with effective time $t_\mathrm{rev}$ and optional compromise time $t_\mathrm{comp}$) and a versioned policy store, and writes an event registry, a content-addressed raw store, a quarantine archive and an evidence log whose signed checkpoints are copied to a witness. Downstream processing derives event-time windows and a current-value table with versioned lineage; an auditor re-verifies the evidence offline. The verifier clock is trusted for horizons and taint; device time stamps are event time, never proof of time.

\begin{figure}[!htbp]
\begin{adjustwidth}{-\extralength}{0cm}
\centering
\includegraphics{figures/fig_architecture.pdf}
\end{adjustwidth}
\caption{Edge-to-cloud data path, trust boundaries and attacker positions AP1--AP6 (Table~\ref{tab:threat}). Solid arrows carry data, dashed arrows control or evidence; $k_o$, $\sigma_o$ and $k_w$, $\sigma_w$ are the origin and relay keys and signatures (Section~\ref{sec:mech}). PEP, policy enforcement point; PDP, policy decision point.\label{fig:arch}}
\end{figure}

Protected assets are observations with their identities, timestamps, sequence information and provenance links, committed raw data, policy history, credentials, lineage, derived data and audit evidence. The trusted computing base comprises the cryptographic primitives, the identity authority's revocation path, the verifier at decision time (including its clock and evidence key, except in insider class T5), policy-history integrity and at least one witness. The device-to-gateway link is assumed not to be attacker-writable (an authenticated fieldbus or device session); Section~\ref{sec:limitations} states what happens otherwise. Table~\ref{tab:threat} lists attacker positions and modelled failures, several of which are indistinguishable by their observable effects. Excluded are breaks of Ed25519 or SHA-256, compromise of the identity authority with all witnesses, process or sensor manipulation before signing, safety and actuation paths, supply-chain and side-channel attacks, confidentiality and flooding attacks.

\begin{table}[!htbp]
\caption{Attacker positions, modelled failures and observational equivalence.\label{tab:threat}}
\begin{adjustwidth}{-\extralength}{0cm}
\footnotesize
\renewcommand\tabularxcolumn[1]{p{#1}}% top-aligned rows
\begin{tabularx}{\fulllength}{l>{\raggedright\arraybackslash}X>{\raggedright\arraybackslash}X}
\toprule
\textbf{Id} & \textbf{Position and capability} & \textbf{Scenarios}\\
\midrule
AP1 & Broker/on-path: modify without keys, inject, drop, delay, replay; claim any publisher identity towards subscribers & E2, E3, E6, E15\\
AP2 & Unknown producer with unregistered credentials & E4\\
AP3 & Stolen device key, injecting at the broker & E5, E14\\
AP4 & Compromised edge workload: signs relay and legacy statements; suppresses, replays, forges for other cells & E4, E15\\
AP3+AP4 & Device key and its authorized relay key: equivocation, gap-fill, full impersonation, back/future dating; forgery under fresh identities is not evaluated (Section~\ref{sec:discussion}) & E5, E14, sensitivity analyses\\
AP5 & Compromised device producing wrong measurements with a valid key & outside the detection scope\\
AP6 & Cloud insider without (T1--T4, T6, T7) or with (T5) the verifier evidence key & E16\\
Failures & Redelivery, acknowledgement loss, loss, reordering, long delay, partition with backlog, gateway restart, device reboot, clock freeze/step, sequence-reuse firmware fault, key rotation during disconnection, buffering beyond the admission horizon & E7--E13, sensitivity analyses\\
\midrule
\multicolumn{3}{l}{\textbf{Observationally equivalent pairs} (benign cause / malicious cause $\rightarrow$ \sys\ treatment)}\\
\multicolumn{3}{l}{identical message twice: retry or redelivery / exact replay $\rightarrow$ duplicate; no state change; labelled by context only}\\
\multicolumn{3}{l}{same identity, different content: sequence-reuse fault / equivocation with a stolen key $\rightarrow$ contested; never authoritative}\\
\multicolumn{3}{l}{missing sequence numbers: loss / suppression by a relay $\rightarrow$ open gap, declared lost after the gap horizon}\\
\multicolumn{3}{l}{very old observation: long disconnection / back-dated forgery or stale replay $\rightarrow$ quarantine, release only by review}\\
\bottomrule
\end{tabularx}
\end{adjustwidth}
\end{table}

\section{PBTI Architecture and Formal Properties}\label{sec:arch}

\subsection{Statements and Admission}\label{sec:mech}

A \sys\ message is a JSON object in the RFC~8785 JSON Canonicalization Scheme (JCS)~\cite{rfc8785}, $\{\Sobs, \sigma_o, \Srel, \sigma_w\}$ (fields in Table~\Sref{tab:S-fields}). The origin signer signs the observation statement with Ed25519~\cite{rfc8032}, $\sigma_o=\mathrm{Sign}(k_o, D_\mathrm{obs}\,\|\,\mathrm{JCS}(\Sobs))$, and the edge workload signs the relay statement, $\sigma_w=\mathrm{Sign}(k_w, D_\mathrm{rel}\,\|\,\mathrm{JCS}(\Srel))$, with domain-separation prefixes $D$. The relay statement contains $\mathit{odig}=H(\mathrm{JCS}(\Sobs)\,\|\,\sigma_o)$ and therefore binds exactly one signed observation. The event identity $e=(\mathit{did},\mathit{ep},\mathit{seq})$ is assigned by the producer and signed; the content digest is $\mathit{cd}=H(\mathrm{JCS}(\Sobs))$ with $H$ = SHA-256, and each observation names its predecessor's digest in \textsf{prev}. Canonicalization and deterministic signatures make benign re-encodings and retransmissions byte-identical; for legacy devices the edge workload assigns $(\mathit{ep},\mathit{seq},\mathit{prev})$ in an epoch namespace unique to that workload.

Algorithm~\ref{alg:admit} lists the admission checks in order; the first failing check determines one of five dispositions (ACCEPT, DUPLICATE, CONFLICT, QUARANTINE, REJECT), and every receipt produces exactly one evidence record. The policy decision point evaluates a versioned, schema-checked document that assigns devices to assets and schemas, states whether a device must sign its observations or may be gateway-attested, and lists the devices each workload may relay; the decision is a pure function of the policy version in force at verifier time $t$ and the statements' identities, so an auditor can re-check it.

\begin{algorithm}[H]
\caption{\sys\ admission of one received message $m$ at verifier time $t$}\label{alg:admit}
\tcp{a reviewed release re-runs this algorithm with only the case's quarantine reason waived}
\footnotesize
$(\Sobs,\sigma_o,\Srel,\sigma_w)\leftarrow$ parse$(m)$ \tcp*{else REJECT malformed}
$K_o\leftarrow$ key$(\Sobs.\mathsf{kid})$, $K_w\leftarrow$ key$(\Srel.\mathsf{wkid})$; require $K_w$ subject $=\Srel.\mathsf{wid}$, and $K_o$ subject $=\Sobs.\mathsf{did}$ (device) or $=\Srel.\mathsf{wid}$ (legacy)\;
verify $\sigma_w$; check $\Srel.\mathsf{odig}=H(\mathrm{JCS}(\Sobs)\|\sigma_o)$; verify $\sigma_o$ \tcp*{else REJECT integrity}
\ForEach{$K\in\{K_o,K_w\}$}{
 \lIf{$K$ revoked at $t$, unless a REVOKED\_BUFFERED waiver meets the release-evidence policy}{\Return QUARANTINE if $\Srel.\mathsf{mode}=$ buffered (policy) else REJECT}
 \lIf{$t>K.\mathit{exp}+\mathit{grace}$}{\Return QUARANTINE expired} \tcp*{within grace: annotate}
}
$(\mathit{permit},\mathit{pv})\leftarrow$ PDP$(\mathsf{did},\mathsf{aid},\mathsf{wid},\mathsf{sch},\mathit{origin},t)$; \lIf{$\neg\mathit{permit}$}{\Return REJECT unauthorized}
\lIf{$\mathsf{ts}>t+\mathit{skew}$ \textbf{or} $t-\mathsf{ts}>h_\mathrm{max}$}{\Return QUARANTINE future/stale}
$s\leftarrow$ slot$(e)$; \lIf{$\mathit{cd}\in s.\mathit{versions}$}{\Return DUPLICATE (RETRY if same relay and within window, else SUSPECTED\_REPLAY)}
\lIf{$e$ in a closed gap, stale epoch, or sequence jump}{\Return QUARANTINE}
add $\mathit{cd}$ to $s.\mathit{versions}$; \lIf{$\mathsf{seq}>0$ \textbf{and} $\Sobs.\mathsf{prev}\neq\bot$}{add $\Sobs.\mathsf{prev}$ to slot$(\mathsf{did},\mathsf{ep},\mathsf{seq}-1).\mathit{implied}$}
remove every slot that changed from COMMITTED to CONTESTED (own or predecessor) from derived data\;
\lIf{$|s.\mathit{versions}\cup s.\mathit{implied}|=1$}{commit raw $m$ under $H(m)$; \Return ACCEPT}
\Return CONFLICT (the version is kept as evidence)\;
\end{algorithm}

\subsection{Rule L and Its Locality}\label{sec:rulel}

For each identity $e$, let $V(e)$ be the set of admitted content digests and $I(e)$ the digests asserted for $e$ by admitted successors' \textsf{prev} links. A slot with an admitted version is COMMITTED when $V(e)\cup I(e)$ contains exactly one distinct digest and CONTESTED when it contains at least two. Without an admitted version it is ABSENT, even if an expected digest is known, or LOST if its gap was closed after the gap horizon; an authorized release can later admit a version into a LOST slot (Table~\Sref{tab:S-transitions}; slot transitions in Figure~\Sref{fig:S-state}). Because versions and implied digests accumulate as sets, the classification is independent of delivery order for a fixed admitted evidence set, a contested identity never becomes committed again, and a gap-fill forgery is contested once a contradictory genuine version or successor link is admitted. Identical contents are idempotent, so the RETRY/SUSPECTED\_REPLAY annotation (same relay, within 30~s) has no authority. Freshness, revocation, policy changes, sequence limits and gap closure can change \emph{which} evidence is admitted, so the pipeline as a whole is not claimed to be invariant under arbitrary arrival schedules.

\begin{Proposition}[Locality of Rule L]\label{prop:locality}
Let $(V,I)$ be an evidence state, $\tau$ the taint flags and $\lambda$ the gap-closure flags, and let $S$ be a sequence of $k$ statements admitted by \sys\ in its main configuration, with no lifecycle transition (gap closure, revocation processing, or a reviewer decision other than the admission it re-runs) in between. Let $A(S)$ be the set of the identities of the statements in $S$ and, for each statement with a positive sequence number and a non-null \textsf{prev}, its predecessor identity. Then $V$ and $I$ change only on $A(S)$; $|A(S)|\le 2k$, with equality exactly when every statement contributes a predecessor and these $2k$ identities are pairwise distinct; and the Rule L class (including LOST) and the authority of every identity outside $A(S)$ are unchanged.
\end{Proposition}
\noindent\emph{Proof sketch.} Admitting $s$ adds $\mathit{cd}(s)$ to $V$ only at $e(s)$ and, if its sequence number is positive and $\mathsf{prev}(s)\neq\bot$, adds $\mathsf{prev}(s)$ to $I$ only at the predecessor of $e(s)$ (Algorithm~\ref{alg:admit}); by induction over $S$, $V$ and $I$ are unchanged outside $A(S)$. The class of an identity $x$ is a function of $V(x)$, $I(x)$ and $\lambda(x)$, and its authority in addition of $\tau(x)$; $\lambda$ is written only by gap closure and $\tau$ only by revocation processing, both excluded. $A(S)$ is the union of at most $k$ statement identities and at most $k$ predecessors. $\square$

The argument is elementary; it was cross-checked by an implementation-independent model and by bounded enumeration on the verifier (Section~\ref{sec:formal}). The bound is attained, for example, by $k$ equivocations at non-adjacent sequence numbers whose \textsf{prev} contradicts a committed predecessor; because admission first advances time-driven lifecycle processing, the hypothesis is a condition on the schedule. B5 (Section~\ref{sec:pipelines}), by contrast, removes every identity of the feed from the fork point onward until review. The proposition bounds the state changes caused by \emph{admitted} contradictory evidence; it does not establish that a contradiction is ever admitted (a suppressed one never is) and excludes lifecycle transitions, admission decisions that earlier admissions change, downstream fan-out, selective holds and replicated verifiers (Section~\Sref{sec:S-protocol}).

\subsection{Credential Lifecycle, Review and Release}\label{sec:lifecycle}

\textbf{Credential lifecycle and retroactive containment.} Expired but unrevoked keys are accepted within a rotation grace period intended for buffered data. After a revocation is observed, ordinary admission under the revoked origin or relay key is refused or, for buffered statements, quarantined (REVOKED\_BUFFERED). If the revocation supplies a compromise time $t_\mathrm{comp}$, committed slots whose \emph{origin} key is revoked and whose verifier commit time is at or after $t_\mathrm{comp}$ are TAINTED: raw bytes are kept, the observation leaves derived processing and affected windows are recomputed. The selector uses commit time, not proof of creation time, so it can exclude legitimate buffered data; a revised record can only move $t_\mathrm{comp}$ earlier. Device-signed observations are not tainted when only their relaying workload was compromised, since the relay could not alter them.

\textbf{Quarantine review and release.} Each quarantined receipt opens a case (QUARANTINED $\rightarrow$ REVIEW\_PENDING $\rightarrow$ RELEASED or REJECTED). A reviewer listed in the policy signs a domain-separated decision bound to the case, the release reason and a snapshot of the evidence examined; the verifier checks the reviewer's authorization and revocation status, the case state and snapshot freshness. A release re-runs admission with only the original quarantine reason waived, so every other check, including Rule L, still applies, and every decision is a signed evidence record that the auditor replays. A REVOKED\_BUFFERED waiver additionally requires a release-evidence condition: \emph{none} (no release under a revoked key), \emph{discretionary} (the observation's own time stamp precedes $t_\mathrm{comp}$; used in every scenario campaign) or \emph{witness} (a token from a time-stamping witness outside the verifier for the exact statement digest, dated before the earliest plausible compromise time). The discretionary condition is a policy condition, not proof of pre-compromise existence (Section~\ref{sec:res-release}). The automated reviewer of the experiments releases a case only if hash links corroborate it: the slot holds that single digest, and the predecessor is authoritative with the digest named by \textsf{prev}, or an admitted successor's \textsf{prev} names it. This is a consistency check under the signers' keys, not independent evidence.

\subsection{Evidence, Lineage and State Model}\label{sec:state}

Accepted messages are stored write-once under $H(m)$, other receipts in the quarantine archive. Evidence records are leaves of RFC~9162 Merkle trees~\cite{rfc9162} over batches of up to 512 records (at least every 5~s); each batch root is signed, chained and copied to the witness. Each derived version stores its sorted input references, their Merkle root and its cause (new, late, contest, taint or release); earlier versions are retained, and the lineage maps onto W3C PROV~\cite{w3c_prov_dm}. The state model uses four orthogonal machines (receipt disposition, event-identity slot, credential lifecycle, versioned derived records) plus a review case per quarantined receipt (Figure~\Sref{fig:S-state}, Table~\Sref{tab:S-transitions}); ten invariants and their oracles are in Table~\Sref{tab:S-invariants}.

\subsection{Verification of the Formal Properties}\label{sec:formal}

The signature bindings rely on the assumed security of Ed25519, SHA-256 and domain separation; there is no computational proof of the complete system. The symbolic model was \VFormalSymbolic\ (ProVerif~\cite{blanchet2016proverif}); these correspondences hold under the model's assumptions. That the implementation realizes the update rule of Proposition~\ref{prop:locality} rests on automated code inspection and bounded checks of an implementation-independent model (\WAfourMcases\ cases), the unmodified verifier with real signatures (\WAfourLexhaustive\ exhaustive and \WAfourLrandom\ sampled batches) and all combinations of up to four lifecycle operations (\WAfourCexhaustive\ exhaustive, \WAfourCrandom\ sampled longer sequences) under nine invariants (\WAfourAllPass; Table~\Sref{tab:S-bounded}). They found one defect, fixed before the confirmatory runs: the taint selector ignored a revised revocation record with an earlier compromise time (no scenario revokes a key twice); an earlier bounded checker found an unrecorded exception for numbers outside the I-JSON range, which all pipelines now reject with an evidence record. Bounded results do not extend to unbounded feeds, arbitrary schedules or replicated verifiers.
